\documentclass[12pt]{beamer}
\usepackage{amsmath, amsfonts, amssymb, amsthm}
\usepackage{graphicx}
\usepackage[brazilian]{babel}
\usepackage[utf8]{inputenc}
\usepackage{verbatim}
\usepackage{hyperref}
\usetheme{Madrid}

\title{Mini-Curso de \LaTeX\ \\ Aula 04 - Matemática}
\author{Fábio Meneghetti \and Pedro Caetano}
\date{5 de abril de 2016}

\begin{document}
\begin{frame}
  \titlepage
\end{frame}

\begin{frame}{Licença}
  Esta apresentação está licenciada com uma Licença Creative Commons Atribuição-CompartilhaIgual 4.0 Internacional.
  \begin{center}
    \includegraphics[scale=0.3]{../license.png}
  \end{center}
\end{frame}

\begin{frame}
  \tableofcontents
\end{frame}

\begin{frame}{Introdução}
  \section{Sobre o modo matemático}
  Para se adicionar equações e fórmulas matemáticas no latex, existe um ambiente dedicado a isso, com comandos específicos que só funcionam quando você está nesse ambiente. Além disso, todas as letras digitadas dentro desse ambiente serão colocadas em itálico, que é o padrão para variáveis matemáticas.
  \\[1cm]
  Mas existem várias formas de se entrar no ambiente matemático, e elas produzem resultados diferentes no documento. Vamos mostrá-las agora.
\end{frame}

\begin{frame}[fragile]{Entrando no modo matemático}
  \section{Entrando no modo matemático}
  \begin{enumerate}
    \item \textbf{O modo em linha (\textit{inline}):} Serve para adicionar matemática no meio de um texto normal. Para entrar nesse modo, basta colocar o código entre cifrões \$.

    \item \textbf{O modo display:} Serve para colocar uma equação matemática em linha própria, centralizada, dando certo destaque. Para entrar nesse modo, basta colocar o código entre \verb+\[+ e \verb+\]+.

    \item \textbf{Equações:} O resultado da equação é igual ao do modo display, porém ela se torna numerada. Desse forma, pode-se fazer referências a equações. Para criar uma equação, basta criar um ambiente equation (com \verb+\begin{equation}+ e \verb+\end{equation}+).

  \end{enumerate}
\end{frame}

\begin{frame}[fragile]
  Antes de fazer um exemplo, vamos aprender alguns comandos simples para poder brincar:
  \begin{itemize}
    \item Para se fazer o ``elevado'', basta utilizar o chapéu \verb+^+. \textbf{Ex:} \verb+x^2+ gera $x^2$
    \item Um valor subescrito é análogo, mas utilizando o \verb+_+. \textbf{Ex:} \verb-x_1 + x_2^3- gera $x_1 + x_2^3$
    \item Para adicionar mais de um só dígito no subescrito ou no elevado, basta ``encapsular'' tudo com chaves. \textbf{Ex:} \verb-e^{2 \pi i}- gera $e^{2 \pi i}$
  \end{itemize}
\end{frame}

\begin{frame}
  Vamos fazer um exemplo! (Exemplo 1)
\end{frame}

\begin{frame}{Símbolos}
  \section{Símbolos}

  A quantidade de símbolos disponíveis para se utilizar no modo matemático é extremamente grande, e não vamos gastar o tempo para mostrando um por um. Mas vamos passar alguns links de páginas com tabelas e listas de símbolos, para que você possa procurar aqueles que você deseja usar. Algumas páginas que julgamos serem boas são:

  \begin{itemize}
    \item \url{http://latex.wikia.com/wiki/List_of_LaTeX_symbols}
    \item \url{https://oeis.org/wiki/List_of_LaTeX_mathematical_symbols}
  \end{itemize}
\end{frame}

\begin{frame}[fragile]{Frações, integrais e somatórios}
  \section{Frações, integrais e somatórios}
  \begin{itemize}
    \item Para adicionar frações, o comando é \verb+\frac{numerador}{denominador}+
    \item O símbolo de integral é \verb+\int+. Ao adicionar valores no subescrito e no elevado da integral, esses valores são posicionados como os extremos de integração. \textbf{Ex:} \verb+\int_0^1 f(x) dx+ gera
    \[
    \int_{0}^1 f(x) dx
    \]
    \item \textbf{Exercício:} O símbolo de somatório é \verb+\sum+. Tente escrever no \LaTeX\ o seguinte somatório:
    \[
      \sum_{k=1}^n x^k
    \]
  \end{itemize}
\end{frame}

\begin{frame}[fragile]{O ambiente align}
  \section{Align}
  Às vezes você precisa escrever várias equações, uma em cima da outra, mas elas possuem cada uma um tamanho diferente, de tal forma que elas não vão ficar alinhadas. Uma solução é o ambiente matemático align, que alinha suas equações.
  \\[5px]

  Para alinhas as equações dentro de um ambiente align, basta separar o lado esquerdo e o direito com um \verb+&+, e usar um \verb+\\+ para terminar a linha.
  \\[5px]

  Dá também para usar vários \verb+&+ para fazer vários alinhamentos. Vamos ver o Exemplo 2!
\end{frame}

\begin{frame}
  \begin{center}
    \large Obrigado!\\
    :)
  \end{center}
\end{frame}

\end{document}
